% Collegamenti tra il capitolo di sicurezza e le tracce passate.
% Incluso da main.tex subito dopo il riassunto.

\section*{Collegamenti con le Tracce} % Usiamo * per non numerare la sezione nell'indice
\addcontentsline{toc}{section}{Collegamenti con le Tracce: Sicurezza}

Nelle tracce raccolte nella cartella \texttt{TraccePassate} (2016--2026) \textbf{non compare un tema dedicato alla sicurezza informatica}. Il capitolo è quindi preventivo. I punti in cui la sicurezza entra comunque in una traccia sono:

\subsection*{Prima sessione 2026 -- A3-INF-4 (prova pratica: cinema multisala)}
Il sistema ha utenti registrati, un'interfaccia Web, API e un database. Un'ottima risposta aggiunge una breve nota di sicurezza: autenticazione, controllo dei ruoli, protezione dagli accessi alle prenotazioni altrui, query parametriche, HTTPS. \textbf{Dove}: sezione \emph{Esempio: Messa in Sicurezza del Sistema di Prenotazione}.

\subsection*{Prima sessione 2026 -- A2-INF-4 (protocolli)}
Tra i protocolli dello stack TCP/IP ``impiegati nei moderni sistemi distribuiti e nelle applicazioni Internet'' rientrano TLS e HTTPS. \textbf{Dove}: sezioni \emph{Infrastruttura a Chiave Pubblica} e \emph{Crittografia Asimmetrica}, e la sezione \emph{Sicurezza di Rete: TLS e HTTPS} del capitolo \emph{Sistemi Operativi e Reti: Approfondimenti}.

\subsection*{Seconda sessione 2024 -- Informatica \#2 (applicazioni distribuite enterprise)}
Tra le ``best practice'' di cloud e architetture distribuite rientrano autenticazione, cifratura, segmentazione e gestione dei segreti. \textbf{Dove}: sezioni \emph{Sicurezza di Rete} e \emph{Sicurezza nel ciclo di sviluppo}.
